\chapter{Frontend}
\label{chap:frontend}

\section{Stack et organisation}

Le frontend (\repo{logiflow-frontend}) est une SPA \textbf{Angular 22} : composants standalone, signals, zoneless, formulaires signal, \texttt{httpResource}. L'UI s'appuie sur \textbf{Tailwind CSS 4} et \textbf{Zard UI}. Cartes Leaflet et graphiques ECharts complètent le tableau de bord.

Les fonctionnalités sont organisées par dossiers métier (\texttt{commandes/}, \texttt{voyages/}, \texttt{maintenance/}, \ldots) avec routes lazy sous un \textbf{SignedInShell}.

\section{App Shell et navigation}

L'\textbf{App Shell} fixe la chrome : sidebar de navigation primaire, zone de contenu scrollable, barre utilitaire (fil d'Ariane, déclencheur Copilote IA, badge File du jour). L'identité et la déconnexion vivent dans le \textbf{pied de sidebar}, pas dans un menu avatar.

La navigation est dérivée d'une table de destinations filtrée par rôles (\texttt{destinationsForRoles}). Un utilisateur ne voit que les modules qu'il a le droit d'ouvrir.

\section{Contrôle d'accès (RBAC)}

Chaque route porte \texttt{data.roles}. Le garde \texttt{roleGuard} redirige vers
\texttt{/403} (Accès refusé) si le rôle de session ne correspond pas. Un utilisateur
authentifié mais non autorisé voit le 403 ; un utilisateur sans session est renvoyé
vers la connexion.

Le tableau~\ref{tab:rbac-ui} résume, pour quelques zones de l'UI, quels rôles métier
y ont typiquement accès (claims portés par le JWT Keycloak).

\begin{table}[H]
\centering
\caption{Zones de l'interface et rôles métier associés (exemples).}
\label{tab:rbac-ui}
\small
\begin{tabularx}{\linewidth}{@{}>{\bfseries\raggedright\arraybackslash}p{3.4cm} >{\raggedright\arraybackslash}X@{}}
\toprule
Zone UI & Rôles autorisés (libellé et claim JWT) \\
\midrule
Commandes / clients &
  Commercial (\texttt{COMMERCIAL}) ;\quad
  Administrateur (\texttt{ADMINISTRATEUR}) \\[0.65em]
Dossiers / voyages &
  Exploitant (\texttt{EXPLOITANT}) ;\quad
  Responsable d'exploitation (\texttt{RESPONSABLE\_EXPLOITATION}) \\[0.65em]
Maintenance &
  Atelier (\texttt{ATELIER}) ;\quad
  Responsable d'exploitation (\texttt{RESPONSABLE\_EXPLOITATION}) \\[0.65em]
Carburant (saisie) &
  Chauffeur (\texttt{CHAUFFEUR}) ;\quad
  plus validateurs côté exploitation \\[0.65em]
Utilisateurs &
  Administrateur (\texttt{ADMINISTRATEUR}) \\
\bottomrule
\end{tabularx}
\end{table}

\section{Authentification}

Le chemin nominal --- en local comme en cloud --- est \textbf{Keycloak} via
\texttt{angular-auth-oidc-client} (Authorization Code + \textbf{PKCE}). Le frontend est un
client public OIDC : après connexion IdP, il joint l'access token sur chaque appel
\texttt{/api/}. Spring agit uniquement comme \emph{resource server} (validation JWT, rôles) ;
il n'émet pas de jetons et ne stocke pas les mots de passe
(voir chapitre~\ref{chap:securite}).

En développement comme en production, le frontend s'authentifie auprès de Keycloak
(OIDC) et le backend (profil \texttt{dev} ou \texttt{prod}) exige un JWT valide sur
l'API métier.

\section{Écrans clés}

Les principaux écrans sont illustrés au chapitre~\ref{chap:captures} :

\begin{itemize}[leftmargin=*]
  \item \textbf{Tableau de bord} : Aperçu (compteurs / breakdowns de statut) puis File du jour.
  \item \textbf{Listes métier} : recherche, filtres, empty states avec CTA unique.
  \item \textbf{Fiches} : détail véhicule, voyage, dossier, avec documents et scores.
  \item \textbf{Nouveau voyage} : sélection de dossiers, calcul d'itinéraire IA, affectations.
  \item \textbf{Copilote} : panneau latéral conversationnel (pas une route dédiée en v1).
  \item \textbf{Palette de commandes} : modules, actions « Nouveau \ldots », références.
\end{itemize}
